\documentclass[11pt]{article}
\usepackage[margin=1in]{geometry}
\usepackage{amsmath,amssymb,amsthm,mathtools}
\usepackage[T1]{fontenc}
\usepackage{lmodern,microtype}
\usepackage[colorlinks=true,allcolors=blue]{hyperref}
\newtheorem{theorem}{Theorem}
\newtheorem{lemma}[theorem]{Lemma}
\newtheorem{proposition}[theorem]{Proposition}
\newcommand{\tr}{\operatorname{Tr}}
\newcommand{\id}{\operatorname{id}}
\newcommand{\norm}[1]{\left\lVert#1\right\rVert}
\newcommand{\ket}[1]{|#1\rangle}
\newcommand{\bra}[1]{\langle#1|}
\title{Operational recovery bounds for an expectation-dependent quantum amplifier}
\author{Daniel Fredriksen\\Quantyra Inc.\\\small\href{https://orcid.org/0009-0002-9123-7435}{ORCID: 0009-0002-9123-7435}}
\date{November 16, 2026\\\small Manuscript candidate 0.1.0; not formally certified}
\begin{document}
\maketitle
\begin{abstract}
We study a hypothetical quantum amplifier whose exchange coupling depends on a source population. For a finite code, we relate correctability of amplitude damping to the non-scalar parts of the compressed source and exchange operators. The bounds use worst-case entanglement recovery error and remain independent of code dimension. A rare-event refinement of the standard information--disturbance relation makes the relevant scale the recovery error divided by damping strength. For trajectories confined to a common invariant moving code, the full nonlinear generator then yields an explicit bound on distinguishability gain in terms of total coupling-time. A finite code and recovery show that the square-root recovery-to-source-compression exponent cannot generally be improved. This does not establish sharpness of the dynamical bound, exclude unprotected excursions or restricted-state recovery, or supply a physical nonlinear device. The contribution is a quantitative application of established approximate-error-correction tools to the stated nonlinear model.
\end{abstract}

\section{Question, baseline and scope}
Ideal nonlinear state discrimination is an established subject~\cite{ChildsYoung}. The existence of an assumed amplifying equation leaves open whether its source can remain useful under error correction. In linear quantum metrology, the relation between the signal and the noise-operator span is already central~\cite{ZhouEtAl,ZhouJiang}. We do not claim that organizing principle as new or import a metrological sufficiency theorem into nonlinear computation.

Our question is narrower. For the model below, how does an \emph{operational recovery guarantee} bound the nonlinear action remaining inside a protected code? An estimate in terms of a non-scalar operator-compression defect alone does not answer this question: that defect is not a physical error probability. We supply an explicit conversion for finite amplitude damping, then integrate the resulting generator bound. The channel theorem of B\'eny and Oreshkov~\cite{BenyOreshkov} is an essential imported result. Our rare-event estimate and finite Kraus identities determine the damping dependence of its application.

All spaces are finite dimensional. The active qubits are a source $S$ and detector $B$, with basis $gg,ge,eg,ee$. Arbitrary additional physical registers are permitted and suppressed in tensor notation. Set
\begin{equation}
 P=\ket g\bra g_S\otimes I_B,\quad N=N_S+N_B,\quad
 J=\ket{eg}\bra{ge}+\ket{ge}\bra{eg},\quad N_S=I-P.
\end{equation}
For a complete normalized pure state $\Psi$, including retained records and passive references, put $p=\langle P\rangle_\Psi$, $j=\langle J\rangle_\Psi$. The native law is
\begin{equation}
 H(\Psi,t)=\omega N+g(t)\sqrt p\,J+
 \frac{g(t)j}{2\sqrt p}(P-pI),\qquad p>0.\label{eq:native}
\end{equation}
Here $\omega$ and $g(t)$ are real. The nonlinear vector field is defined to be zero at $p=0$. The last term is required by the gradient of the homogeneous functional
\begin{equation}
 \mathcal E(\Psi)=\omega\langle\Psi,N\Psi\rangle+
 g\sqrt{\frac{\langle\Psi,P\Psi\rangle}{\langle\Psi,\Psi\rangle}}
 \langle\Psi,J\Psi\rangle.\label{eq:functional}
\end{equation}
We retain this term throughout. In particular, replacing the unknown population by a classically supplied number would be a different operation. No physical realization, causal completion or operational energy apparatus for~\eqref{eq:native} is supplied here.

The trajectory theorem assumes that the compared clean trajectories remain in the range of a common differentiable isometry $V(t)$ during native evolution. Ordinary controls and changes of frame are independent of the unknown input. This is a stronger hypothesis than the requirements of an arbitrary successful nonlinear computer. Our recovery test assesses a code against a specified ordinary noise channel with native evolution switched off; it is not a model of an arbitrary simultaneous noisy nonlinear flow.

\section{Recovery error and a rare-event estimate}
Write $D(\rho,\sigma)=\tfrac12\norm{\rho-\sigma}_1$ and use root fidelity
\begin{equation}
 f(\rho,\sigma)=\norm{\sqrt\rho\sqrt\sigma}_1,\qquad
 F(\mathcal A,\mathcal B)=\min_\rho
 f\bigl((\mathcal A\otimes\id)(\psi_\rho),
        (\mathcal B\otimes\id)(\psi_\rho)\bigr),
\end{equation}
where $\psi_\rho$ is a purification. For an encoded noise channel $\mathcal N_V$, suppose an ordinary recovery $\mathcal R$ satisfies
\begin{equation}
 F(\mathcal R\mathcal N_V,\id)^2\geq1-\eta,
 \qquad 0\leq\eta\leq1.\label{eq:recovery}
\end{equation}
Thus $\eta$ is a worst-case entanglement-infidelity bound, not the square root of that error. A half-diamond recovery-error bound $\tfrac12\norm{\mathcal R\mathcal N_V-\id}_\diamond\leq\eta$ suffices for~\eqref{eq:recovery}: on each purified input the target is pure and its probability decreases by at most $\eta$.

B\'eny--Oreshkov's channel fidelity duality (Theorem~1, equation~(3) of Ref.~\cite{BenyOreshkov}), specialized to the identity target, implies that some environment state $\sigma$ satisfies
\begin{equation}
 F(\widehat{\mathcal N}_V,\mathcal S_\sigma)\geq\sqrt{1-\eta},
 \qquad \mathcal S_\sigma(\rho)=\sigma\tr\rho.\label{eq:duality}
\end{equation}
Here the hat denotes a complementary channel. This is an existence statement. It does not provide a low-cost algorithm for finding $V$, $\mathcal R$ or $\sigma$. In particular, no semidefinite-program cost is hidden in a computational claim.

\begin{lemma}[Rare-event observable bound]\label{lem:rare}
Let $\Pi$ be a projector and $O$ Hermitian with $O=\Pi O\Pi$ and $\norm O\leq1$. Suppose $f(\rho,\sigma)\geq\sqrt{1-\eta}$ and $\tr(\Pi\rho)\leq w$. Set
\begin{equation}
 b(\eta)=\sqrt{2(1-\sqrt{1-\eta})}\leq\sqrt{2\eta}.
\end{equation}
Then
\begin{equation}
 |\tr O(\rho-\sigma)|\leq2b(\eta)\sqrt w+b(\eta)^2.\label{eq:rare}
\end{equation}
\end{lemma}
\begin{proof}
By the purification characterization of fidelity, choose unit purifications $x,y$ of $\rho,\sigma$, with the phase of $y$ chosen so that $\langle x,y\rangle=f(\rho,\sigma)\geq0$. Then $\norm{x-y}^2=2-2f(\rho,\sigma)\leq b^2$, where $b=b(\eta)$. Extend $O,\Pi$ by the identity on the purifying space. Then $\norm{\Pi x}\leq\sqrt w$ and $\norm{\Pi y}\leq\sqrt w+b$. Expanding the difference of expectations and applying Cauchy--Schwarz gives
\[
 |\langle x,Ox\rangle-\langle y,Oy\rangle|
 \leq\norm{x-y}\bigl(\norm{\Pi x}+\norm{\Pi y}\bigr)
 \leq b(2\sqrt w+b).
\]
Finally $1-\sqrt{1-\eta}\leq\eta$. The argument also covers $w=0$ and $\eta=0$.
\end{proof}
The purification and fidelity facts used here are standard~\cite{Watrous}. The refinement over a plain trace-distance estimate is the factor $\sqrt w$ when the relevant environmental event is rare.

\section{From finite damping recovery to source compression}
For $0<\gamma\leq1$, let one-qubit amplitude damping have Kraus operators
\begin{equation}
 A_0=\ket g\bra g+\sqrt{1-\gamma}\ket e\bra e,
 \qquad A_1=\sqrt\gamma\ket g\bra e.
\end{equation}
On $SB$, put $K_{ij}=A_i\otimes A_j$, $i,j\in\{0,1\}$. The channel acts as the identity on all other physical registers. Compose it with the encoding $\rho\mapsto V\rho V^\dagger$ to obtain $\mathcal N_V$. Label the complementary environment basis by $00,10,01,11$, and define
\begin{equation}
 \Pi_S=\ket{10}\bra{10}+\ket{11}\bra{11},\qquad
 T=\ket{10}\bra{01}+\ket{01}\bra{10}.
\end{equation}
Direct Kraus multiplication gives
\begin{align}
 \widehat{\mathcal N}_V^*(\Pi_S)&=\gamma V^\dagger N_S V,\label{eq:pullP}\\
 \widehat{\mathcal N}_V^*(T)&=\gamma V^\dagger J V.\label{eq:pullJ}
\end{align}
For example, $K_{10}^\dagger K_{01}=\gamma\ket{eg}\bra{ge}$, including at $\gamma=1$. The single-jump projector $\Pi_1=\ket{10}\bra{10}+\ket{01}\bra{01}$ obeys
\begin{equation}
 \widehat{\mathcal N}_V^*(\Pi_1)
 =\gamma V^\dagger(N_S+N_B-2\gamma N_SN_B)V.
\end{equation}
Consequently its probability is at most $2\gamma$, while the source-jump probability is at most $\gamma$. These bounds are deliberately conservative when $\gamma$ is large.

\begin{theorem}[Operational compression bound]\label{thm:compression}
For the stated two-qubit damping test with $0<\gamma\leq1$ and identity action on other registers, under~\eqref{eq:recovery} there are real scalars $a,c$ such that
\begin{align}
 \norm{V^\dagger PV-aI}&\leq e_P(\gamma,\eta),\\
 \norm{V^\dagger JV-cI}&\leq e_J(\gamma,\eta),\label{eq:compression}
\end{align}
where, with $b=b(\eta)$,
\begin{align}
 e_P(\gamma,\eta)&=\min\left\{\tfrac12,
                 \frac{2b}{\sqrt\gamma}+\frac{b^2}{\gamma}\right\},\\
 e_J(\gamma,\eta)&=\min\left\{1,
                 \frac{2\sqrt2b}{\sqrt\gamma}+\frac{b^2}{\gamma}\right\}.\label{eq:eps}
\end{align}
The result is independent of logical and spectator dimensions.
\end{theorem}
\begin{proof}
Use the state $\sigma$ from~\eqref{eq:duality}. For every logical pure input, apply Lemma~\ref{lem:rare} to its environment output: first $O=\Pi_S$ with $w=\gamma$, then $O=T$ supported on $\Pi_1$ with $w=2\gamma$. Divide the resulting bounds by $\gamma$ and use~\eqref{eq:pullP}--\eqref{eq:pullJ}. Since the expectations of a Hermitian matrix over all unit vectors determine its operator norm, this proves scalar-approximation bounds with the second arguments of the minima in~\eqref{eq:eps}. The optimal real scalar is the midpoint of the smallest and largest eigenvalues. Its error is half the spectral diameter and is no larger than either obtained bound. Since $0\leq V^\dagger PV\leq I$ and $\norm{V^\dagger JV}\leq1$, the trivial midpoint bounds are $1/2$ and $1$, respectively. Choosing the midpoint scalars proves both minima simultaneously, without computing $\sigma$.
\end{proof}
In particular $\eta=0$ makes both compressed operators scalar. The theorem uses an ordinary finite noise channel, including no-jump and double-jump terms; it does not replace finite damping by an informal single-loss test.

\section{A whole-trajectory distinguishability bound}
Let two complete clean trajectories have initial trace distance $d_0$. A trajectory here means an absolutely continuous solution on a finite concatenation of compact native-on intervals. During each interval the states lie in the range of a common $C^1$ isometry $V(t)$, possibly tensored with passive references. The real coupling $g$ and common Hermitian ordinary generator are integrable, and $\gamma,\eta$ are measurable. Suppose the code at every such time satisfies~\eqref{eq:recovery} for the finite damping test with $0<\gamma(t)\leq1$, $0\leq\eta(t)\leq1$. Ordinary controls and the moving frame are common to both inputs.

Each boundary between native intervals is an input-independent isometry on the complete space: an ordinary channel is implemented by a specified Stinespring dilation, with every environment and record retained through all subsequent native evolution. Each subsequent interval has its own common code containing both complete states. The native operators and damping test act as the identity on all spectators, including these records. Classical feed-forward is represented by a common coherent control on retained records. No register is discarded before the last native interval, and no separately evolved unread nonlinear branches are assumed. The final ordinary readout may discard registers.

For finitely many simultaneous pairs $k$, define the native nonlinear part of the complete generator explicitly by
\begin{equation}
 H_{\mathrm{nl}}(\Psi,t)=\sum_k\left[
 g_k\sqrt{p_k}J_k+\frac{g_kj_k}{2\sqrt{p_k}}(P_k-p_kI)\right],\label{eq:multi}
\end{equation}
with the stated zero-population extension for each summand and a common integrable ordinary Hermitian part. Every $g_k$ is real and integrable; $p_k,j_k$ are expectations on the complete state. For each pair separately assume~\eqref{eq:recovery} for its own two-qubit damping test, identity elsewhere, with measurable parameters $\gamma_k,\eta_k$. The same common code must satisfy all these tests. Pairs may overlap; no commutation between different summands is required.

\begin{theorem}[Gain limited by operational recovery]\label{thm:gain}
Under these hypotheses, a single native pair satisfies
\begin{equation}
 D_{\rm out}\leq\min\left\{1,\ d_0+
 2\int |g(t)|[e_P(\gamma(t),\eta(t))+e_J(\gamma(t),\eta(t))]dt\right\}.\label{eq:gain}
\end{equation}
For model~\eqref{eq:multi}, sum the integrands for each pair with its own $g_k,\gamma_k,\eta_k$.

If $\eta_k(t)/\gamma_k(t)\leq r\leq1$ for every active pair and time, define $A=\int\sum_k|g_k(t)|dt$, with one summand for a single pair. Then
\begin{equation}
 D_{\rm out}\leq\min\{1,d_0+22A\sqrt r\}.\label{eq:simple}
\end{equation}
\end{theorem}
\begin{proof}
Write $V^\dagger PV=aI+E_P$, $V^\dagger JV=cI+E_J$ using midpoint scalars. The ordinary part of the logical generator consists of the compressed ordinary controls, $\omega V^\dagger NV$, and the common moving-frame connection $-iV^\dagger\dot V$. Removing scalar terms from the nonlinear part of~\eqref{eq:native} leaves
\begin{equation}
 R_\Psi=g\sqrt p\,E_J+\frac{gj}{2\sqrt p}E_P.\label{eq:residual}
\end{equation}
Because $PJP=(I-P)J(I-P)=0$ and $\norm J\leq1$,
\begin{equation}
 |j|\leq2\sqrt{p(1-p)},\qquad
 \norm{R_\Psi}\leq|g|(e_J+e_P).\label{eq:bounded}
\end{equation}
This estimate does not divide by a lower bound on $p$. At $p=0$ only the native nonlinear vector field (the $g$-dependent terms) is zero; the ordinary contribution $\omega N\Psi$ is retained. The nonlinear density-operator contribution then has a zero representative satisfying the same bound. To see the continuity of that nonlinear vector contribution, $\norm{P\Psi}=\sqrt p$ and~\eqref{eq:bounded}'s first inequality bound each apparently singular term by a constant times $|g|\sqrt p$. The residual along an actual trajectory is measurable: spectral midpoints and the state expectations are continuous, and the quotient on $p>0$ with the specified value on $p=0$ is measurable. Its norm is bounded by $\tfrac32|g|$, so it is integrable.

For either input, compare the logical state with evolution under the common ordinary generator alone. Variation of constants and
\begin{equation}
 \tfrac12\norm{[R,\rho]}_1\leq\norm R
\end{equation}
give a trace-distance change of at most $\int|g|(e_P+e_J)dt$. This calculation evaluates $R$ on the actual trajectory; it does not assume the nonlinear flow is contractive. The triangle inequality for the two actual inputs supplies the factor two in~\eqref{eq:gain}. The specified common boundary isometries preserve trace distance, and the final ordinary readout contracts it. Apply the comparison anew to the two actual states at each interval start and sum the increments. For~\eqref{eq:multi}, apply Theorem~\ref{thm:compression} to each pair's separate test and sum residual norms; overlap between pairs does not change this triangle bound.

Finally $b/\sqrt\gamma\leq\sqrt{2\eta/\gamma}$ and $b^2/\gamma\leq2\eta/\gamma$. Dropping the caps, the integrand factor obeys
\[
 e_P+e_J\leq(2\sqrt2+4)\sqrt r+4r
 \leq(8+2\sqrt2)\sqrt r<11\sqrt r.
\]
This proves~\eqref{eq:simple}.
\end{proof}

Equal-prior discrimination with success at least $2/3$ requires $D_{\rm out}\geq1/3$. Hence for $d_0\leq1/4$ the uniform bound implies
\begin{equation}
 A\sqrt r\geq\frac1{264}.\label{eq:cost}
\end{equation}
If $r$ is exponentially small, this particular protected-trajectory class needs exponentially large coupling-time $A$. A polynomial bound on the summed native strength then implies exponentially large duration. Conversely, inverse-polynomial $r$ over polynomial coupling-time is not excluded. The bound does not prove that such an algorithm exists.

The damping normalization is essential. If a noise slice has $\gamma\simeq\kappa\Delta t$, the relevant ratio is $\eta/(\kappa\Delta t)$, not $\eta$ alone. Absolute recovery errors going to zero as $\Delta t\to0$ therefore do not automatically suppress the amplifier. This paper does not derive $\eta(\Delta t)$ for a physical control implementation, pay an assumed infinitely fast correction cycle, or replace the common-invariant-code hypothesis by a theorem about arbitrary excursions. Encoding, recovery, native synthesis, control accuracy and classical construction costs remain additional requirements. Only the explicitly defined summed native coupling-time is lower bounded here.

\section{The square-root compression exponent is necessary}
The following construction concerns the recovery-to-compression step, not attainment of the dynamical gain bound. Fix $0<a<1$ and $0<d<\min(a,1-a)$; put $n_0=a+d$, $n_1=a-d$. With a four-dimensional noiseless ancillary register labelled by $(j,f)\in\{0,1\}^2$ and detector $B$ fixed in $g$, encode
\begin{equation}
 V\ket j=\sqrt{1-n_j}\ket{gg}\ket{j,0}
             +\sqrt{n_j}\ket{eg}\ket{j,1}.\label{eq:code}
\end{equation}
The code is an isometry and
\begin{equation}
 \inf_{c\in\mathbb R}\norm{V^\dagger PV-cI}=d.\label{eq:defect}
\end{equation}
Apply the same two-qubit damping channel; the detector has no excitation to lose. For each $j$, the no-jump and source-jump vectors have norms $\sqrt{1-\gamma n_j}$ and $\sqrt{\gamma n_j}$. Their four normalized directions are orthogonal: the jump has source $g$ and ancillary flag $1$, while the no-jump vector has support only on $(g,0)$ and $(e,1)$. An ordinary isometric decoder maps them to $\ket j\ket{\mathrm{no}}$ and $\ket j\ket{\mathrm{jump}}$, respectively. Complete that isometry with extra syndrome dimensions if needed, then trace the syndrome.

The resulting logical channel preserves diagonal entries and multiplies off-diagonal entries by
\begin{equation}
 \lambda=\sqrt{(1-\gamma a)^2-\gamma^2d^2}
                 +\gamma\sqrt{a^2-d^2}.\label{eq:lambda}
\end{equation}
It is the phase-flip channel with error probability $\eta_{\rm rec}=(1-\lambda)/2$. For a purified logical input with reduced state $\rho$, its squared entanglement fidelity is
\begin{equation}
 1-\eta_{\rm rec}+\eta_{\rm rec}(\tr\rho Z)^2.
\end{equation}
Thus its worst-case infidelity is exactly $\eta_{\rm rec}$, attained when $\tr\rho Z=0$. Expanding the two positive square roots at $d=0$, for fixed $a,\gamma$ gives
\begin{equation}
 \eta_{\rm rec}=\frac{\gamma d^2}{4a(1-\gamma a)}+O(d^4).\label{eq:scaling}
\end{equation}
In particular, a universal replacement of the recovery-to-source-defect scale $\sqrt{\eta/\gamma}$ by $o(\sqrt{\eta/\gamma})$ would fail on this family: the achieved error in~\eqref{eq:scaling} is a legitimate recovery guarantee while the defect is exactly $d$. This conclusion remains true if one substitutes the optimal recovery error, which is no larger than the achieved one.

This finite decoder depends only on known code and damping parameters, not an unknown logical state. We make no optimal synthesis-cost claim. Crucially, $J$ sends the excited-source component of~\eqref{eq:code} to a detector-excited vector outside the code. The construction therefore does \emph{not} satisfy the native-invariance premise of Theorem~\ref{thm:gain} and does not demonstrate a working protected nonlinear amplifier. Neither the constant 22 nor the overall dynamical exponent is proved optimal.

\section{Contribution, limitations and verification}
The operational bridge replaces unspecified compression defects by a standard worst-case recovery error, with explicit damping dependence and no code-dimension loss. The channel-fidelity duality is imported, and the signal/noise compatibility principle is established prior work. The new manuscript contribution is the stated combination of the finite damping observables, rare-event estimate, complete nonlinear-generator bound and finite scaling witness. Exhaustive novelty and priority are not certified.

The theorem does not address nonlinear recovery, state families smaller than an unknown logical code, input-dependent classical computation of the answer, arbitrary finite excursions from a protected code, a native interaction outside the tested noise-observable structure, or physical realization. It does not convert a conserved scalar functional into an operational total-energy theorem. Ordinary BPP versus BQP is unchanged.

The proofs are analytic. A fixed exact diagnostic checks finite Kraus identities and the encoding at rational parameter choices; no numerical agreement establishes the asymptotic theorem. The repository preserves the antecedent research notes with their original limitations. Separate proof-adversarial, complexity/resource and non-claims/attribution reports and their source hashes are tracked in the repository; those receipts determine the review status. Such AI reviews do not constitute Lean certification or human peer review.

\paragraph{Assistance and status.} OpenAI Codex assisted with derivation, source comparison, drafting and internal checks; separate Claude, Grok and AGY reviews are recorded as described above. This is a preparation-stage candidate, not a physical proposal or an archival publication. Author funding and competing-interest declarations remain to be finalized for any selected submission destination.

\begin{thebibliography}{9}
\raggedright
\bibitem{BenyOreshkov} C. B\'eny and O. Oreshkov, ``General Conditions for Approximate Quantum Error Correction and Near-Optimal Recovery Channels,'' \emph{Phys. Rev. Lett.} \textbf{104}, 120501 (2010). \href{https://doi.org/10.1103/PhysRevLett.104.120501}{doi:10.1103/PhysRevLett.104.120501}. \href{https://arxiv.org/abs/0907.5391v4}{arXiv:0907.5391v4}.
\bibitem{ZhouEtAl} S. Zhou, M. Zhang, J. Preskill and L. Jiang, ``Achieving the Heisenberg limit in quantum metrology using quantum error correction,'' \emph{Nature Communications} \textbf{9}, 78 (2018). \href{https://doi.org/10.1038/s41467-017-02510-3}{doi:10.1038/s41467-017-02510-3}. \href{https://arxiv.org/abs/1706.02445v2}{arXiv:1706.02445v2}.
\bibitem{ZhouJiang} S. Zhou and L. Jiang, ``Optimal Approximate Quantum Error Correction for Quantum Metrology,'' \emph{Phys. Rev. Research} \textbf{2}, 013235 (2020). \href{https://doi.org/10.1103/PhysRevResearch.2.013235}{doi:10.1103/PhysRevResearch.2.013235}. \href{https://arxiv.org/abs/1910.08472v1}{arXiv:1910.08472v1}.
\bibitem{ChildsYoung} A. M. Childs and J. Young, ``Optimal state discrimination and unstructured search in nonlinear quantum mechanics,'' \emph{Phys. Rev. A} \textbf{93}, 022314 (2016). \href{https://doi.org/10.1103/PhysRevA.93.022314}{doi:10.1103/PhysRevA.93.022314}. \href{https://arxiv.org/abs/1507.06334v2}{arXiv:1507.06334v2}.
\bibitem{Watrous} J. Watrous, \emph{The Theory of Quantum Information}, Cambridge University Press (2018). Author's text: \href{https://cs.uwaterloo.ca/~watrous/TQI/}{cs.uwaterloo.ca/\textasciitilde watrous/TQI/}.
\end{thebibliography}
\end{document}
